\documentclass[11pt,a4paper]{article}
\usepackage{geometry}
\usepackage{url}
\geometry{margin=27mm}

\title{arabtex-granada\\[3pt]\large Package guide}
\author{Francisco M. Garc\'ia}
\date{Version 1.0.0 \quad 15 August 2026}

\begin{document}
\maketitle

\begin{abstract}
\noindent
\texttt{arabtex-granada} is a self-contained adaptation of Klaus Lagally's
ArabTeX. It preserves the historical Arabic typesetting runtime and adds the
transcription convention used by the School of Arabists of Granada. This guide
covers the new package; it is not a replacement for the original ArabTeX manual.
\end{abstract}

\section{Quick start}

Load the package, select Arabic text, then select the Granada transcription:

\begin{verbatim}
\documentclass{article}
\usepackage{arabtex-granada}
\begin{document}
\setarab
\settrans{spanish}
\arabtrue\transtrue\RL{al-`arabiyya}
\end{document}
\end{verbatim}

The package renders transliterated ArabTeX input as Arabic script and, when
transcription is enabled, as a Latin transcription. For a complete example
demonstrating each changed Granada mapping, see
\begin{center}
\path{examples/arabtex-granada-example.tex}
\end{center}

\section{Granada transcription}

\verb|\settrans{spanish}| starts with ArabTeX's ZDMG transcription and changes
only the following input-to-output mappings:

\begin{center}
\begin{tabular}{ll}
\textbf{ArabTeX input} & \textbf{Granada transcription} \\
\hline
\verb|^g| & \^{y} \\
\verb|_h| & j \\
\verb|.g| & g
\end{tabular}
\end{center}

This option changes transcription, not the underlying Arabic input syntax.
For ArabTeX's full input notation and commands, consult the upstream ArabTeX
documentation at \url{https://ctan.org/pkg/arabtex}.

\section{Installation}

The CTAN archive includes the package entry point, private ArabTeX runtime,
font metrics, Type 1 fonts, Metafont sources, and \texttt{arabtex.map}. The
upstream \texttt{arabtex} CTAN package is not a dependency and should not be
loaded alongside this package.

For manual installation, place \path{arabtex-granada.sty} and
\path{vendor/texinput/} in \path{tex/latex/arabtex-granada/}.
Install each bundled font directory under the matching TeX Directory
Structure \path{fonts/} category: metrics in \path{tfm/}, Type 1 fonts in
\path{type1/}, Metafont sources in \path{source/}, and the map in
\path{map/dvips/}. See \path{README.md} for the exact destination paths.
Refresh the filename database and enable \texttt{arabtex.map} using your TeX
distribution's tools. TeX Live and MiKTeX installations normally handle
these steps automatically.

If a local PDF build reports a missing ArabTeX font or font map, check that
the bundled metrics and Type 1 fonts are visible and that
\texttt{arabtex.map} is enabled. LaTeX source parsing alone does not prove
that PDF font embedding succeeded.

\section{Compatibility and provenance}

The adaptation preserves modern LaTeX kernel definitions of
\verb|\begin| and \verb|\end| while loading the historical ArabTeX runtime.
This avoids package-hook failures when another package is loaded at the start
of the document. On older LaTeX formats without the modern hook interface,
ArabTeX's historical environment behavior remains in place.

Klaus Lagally authored the original ArabTeX distribution. Francisco M.
Garc\'ia maintains this adaptation and its Granada convention. The bundled
original files retain their authorship and history; see
\path{vendor/README.md} and \path{vendor/ORIGINAL_README.txt}.
This derived work is distributed under LPPL version 1.3c or later; see
\texttt{lppl.txt}. Problems with this adaptation should be reported to its
maintainer, not to the original ArabTeX author.

\end{document}
